\documentclass[11pt,a4paper]{article}
\usepackage[utf8]{inputenc}
\usepackage{amsmath,amssymb}
\usepackage{geometry}
\usepackage{booktabs}
\usepackage{xcolor}
\usepackage{listings}
\usepackage{hyperref}

\geometry{margin=1in}

\definecolor{codegreen}{rgb}{0,0.6,0}
\definecolor{codegray}{rgb}{0.5,0.5,0.5}
\definecolor{codepurple}{rgb}{0.58,0,0.82}
\definecolor{backcolour}{rgb}{0.96,0.96,0.96}

\lstdefinestyle{mystyle}{
    backgroundcolor=\color{backcolour},   
    commentstyle=\color{codegreen},
    keywordstyle=\color{blue},
    numberstyle=\tiny\color{codegray},
    stringstyle=\color{codepurple},
    basicstyle=\ttfamily\footnotesize,
    breakatwhitespace=false,         
    breaklines=true,                 
    captionpos=b,                    
    keepspaces=true,                 
    numbers=left,                    
    numbersep=5pt,                  
    showspaces=false,                
    showstringspaces=false,
    showtabs=false,                  
    tabsize=2
}
\lstset{style=mystyle}

\title{\textbf{Day 1: Python, Pandas, and Data Preprocessing for Machine Learning}}
\author{Data Science Foundation Series}
\date{\today}

\begin{document}

\maketitle

\tableofcontents
\newpage

\section{Introduction and Architecture}
Data preprocessing forms 80\% of a Machine Learning Engineer's daily workflow. Machine Learning algorithms rely on matrix algebra:
\begin{equation}
\mathbf{\hat{y}} = \sigma(\mathbf{X}\mathbf{W} + \mathbf{b})
\end{equation}
Because algorithms perform linear and non-linear transformations on numeric matrices $\mathbf{X} \in \mathbb{R}^{n \times d}$, raw real-world data cannot be consumed directly. Raw data contains missing values, irregular strings, high cardinality text, and unscaled magnitudes.

\section{Pandas Foundations: Series and DataFrames}
\subsection{Series}
A Series is a one-dimensional labeled homogeneous array. Mathematically, it is a function mapping an index set to value space:
\begin{equation}
f: \mathcal{I} \to \mathcal{V}
\end{equation}

\subsection{DataFrames and Indexing}
A DataFrame is a two-dimensional tabular data structure with heterogeneous column types:
\begin{itemize}
    \item \texttt{.loc}: Label-based indexing (uses column names and row labels).
    \item \texttt{.iloc}: Integer-position indexing (0 to $N-1$).
\end{itemize}

\section{Data Cleaning Masterclass}
\subsection{Duplicate Resolution}
Duplicates artificially inflate model confidence and bias predictions. We use:
\begin{lstlisting}[language=Python]
df = df.drop_duplicates(subset=['Emp_ID'], keep='first')
\end{lstlisting}

\subsection{Missing Data Imputation Strategies}
Let $x_i$ represent the $i$-th sample of feature $X$:
\begin{enumerate}
    \item \textbf{Mean Imputation}: $\mu = \frac{1}{N} \sum_{i=1}^N x_i$ (Sensitive to extreme outliers).
    \item \textbf{Median Imputation}: Robust central tendency, minimizes $L_1$ loss:
    \begin{equation}
    \text{Median}(X) = \arg\min_c \sum_{i=1}^N |x_i - c|
    \end{equation}
    \item \textbf{Categorical Mode}: Most frequent category $c^* = \arg\max_c P(c)$.
\end{enumerate}

\subsection{Outlier Handling via IQR}
Given the first quartile $Q_1$ (25th percentile) and third quartile $Q_3$ (75th percentile), the Interquartile Range (IQR) is:
\begin{equation}
\text{IQR} = Q_3 - Q_1
\end{equation}
Any observation $x$ outside the interval is considered an outlier:
\begin{equation}
x \notin [Q_1 - 1.5 \times \text{IQR}, \; Q_3 + 1.5 \times \text{IQR}]
\end{equation}

\section{Exploratory Data Analysis (EDA)}
The correlation between two continuous random variables $X$ and $Y$ is given by Pearson's coefficient:
\begin{equation}
\rho_{X, Y} = \frac{\text{Cov}(X, Y)}{\sigma_X \sigma_Y} = \frac{\sum_{i=1}^n (x_i - \bar{x})(y_i - \bar{y})}{\sqrt{\sum_{i=1}^n (x_i - \bar{x})^2}\sqrt{\sum_{i=1}^n (y_i - \bar{y})^2}}
\end{equation}

\section{Machine Learning Preprocessing}
\subsection{Feature Encoding}
\textbf{One-Hot Encoding}: Maps a categorical variable with $K$ distinct categories into $K-1$ binary columns to prevent multicollinearity (the dummy variable trap).

\subsection{Train-Test Splitting \& Data Leakage}
\textbf{The Golden Law:} The test dataset $\mathcal{D}_{\text{test}}$ must remain completely invisible during all preprocessing stages. Never compute mean $\mu$ or variance $\sigma^2$ on the full dataset before splitting.

\subsection{Feature Scaling}
\subsubsection{Standardization (Z-Score)}
\begin{equation}
z_i = \frac{x_i - \mu_{\text{train}}}{\sigma_{\text{train}}}
\end{equation}
Transforms distribution to mean 0 and standard deviation 1.

\subsubsection{Normalization (Min-Max Scaling)}
\begin{equation}
x_{i,\text{norm}} = \frac{x_i - x_{\min,\text{train}}}{x_{\max,\text{train}} - x_{\min,\text{train}}}
\end{equation}
Bounds features strictly to $[0, 1]$.

\end{document}
